% Auto-gerado por C8_build_T1_latex.py - NÃO editar à mão
\begin{table}[htbp]
\centering
\caption{Indicadores agregados de fluxo escolar, Brasil, por macroetapa e ano (PNADC longitudinal, 2012--2023).}
\label{tab:t1_brasil}
\begin{threeparttable}
\small
\begin{tabular}{lrrrrrr}
\toprule
Ano & Promoção & Repetência & Evasão & Não-progr. & Abandono$^{a}$ & N \\
    & (\%)     & (\%)        & (\%)   & (\%)        & (\%)             &   \\
\midrule
\multicolumn{7}{l}{\textit{Painel A: EF iniciais (1\textsuperscript{o}--5\textsuperscript{o} ano)}} \\
\midrule
2012 & 65.4 & 32.2 & 1.3 & 33.6 & 1.4 & 31.321 \\
2013 & 66.8 & 30.6 & 1.2 & 32.0 & 1.2 & 28.844 \\
2014 & 71.2 & 26.7 & 1.2 & 28.0 & 1.1 & 28.359 \\
2015 & 73.6 & 25.7 & 0.5 & 26.4 & 0.8 & 20.668 \\
2016 & 73.8 & 25.6 & 0.5 & 26.1 & 0.5 & 30.230 \\
2017 & 74.4 & 24.9 & 0.4 & 25.5 & 0.5 & 28.904 \\
2018 & 75.6 & 23.9 & 0.4 & 24.4 & 0.4 & 28.131 \\
2019 & 76.6 & 22.9 & 0.3 & 23.3 & 0.5 & 26.102 \\
2020$^{\dagger}$ & 57.0 & 42.2 & 0.8 & 43.0 & 0.6 & 15.736 \\
2021$^{\dagger}$ & 65.0 & 34.6 & 0.3 & 34.9 & 0.4 & 21.463 \\
2022 & 56.4 & 43.4 & 0.2 & 43.6 & 0.2 & 22.008 \\
2023 & 56.0 & 43.8 & 0.2 & 44.0 & 0.1 & 22.787 \\
\\[0.3em]
\multicolumn{7}{l}{\textit{Painel B: EF finais (6\textsuperscript{o}--9\textsuperscript{o} ano)}} \\
\midrule
2012 & 59.9 & 32.7 & 3.2 & 36.8 & 3.3 & 23.077 \\
2013 & 62.2 & 30.5 & 3.1 & 34.2 & 3.2 & 21.947 \\
2014 & 65.4 & 28.3 & 2.8 & 31.9 & 3.0 & 22.075 \\
2015 & 66.7 & 28.0 & 2.2 & 31.3 & 2.4 & 16.167 \\
2016 & 67.5 & 27.8 & 1.8 & 30.7 & 2.1 & 24.293 \\
2017 & 69.0 & 26.8 & 1.7 & 29.6 & 2.0 & 23.514 \\
2018 & 71.1 & 24.8 & 1.8 & 27.4 & 1.8 & 22.774 \\
2019 & 71.5 & 24.7 & 1.6 & 27.0 & 1.7 & 21.342 \\
2020$^{\dagger}$ & 53.7 & 44.1 & 1.1 & 45.5 & 0.6 & 13.624 \\
2021$^{\dagger}$ & 63.4 & 34.2 & 1.3 & 35.9 & 0.6 & 18.854 \\
2022 & 54.6 & 43.7 & 0.9 & 44.8 & 0.8 & 18.923 \\
2023 & 54.6 & 44.1 & 0.5 & 44.9 & 0.6 & 19.126 \\
\\[0.3em]
\multicolumn{7}{l}{\textit{Painel C: Ensino M\'edio (1\textsuperscript{o}--3\textsuperscript{o} ano)}} \\
\midrule
2012 & 44.9 & 20.7 & 24.4 & 46.1 & 11.8 & 14.700 \\
2013 & 45.0 & 20.3 & 23.8 & 45.0 & 11.6 & 14.071 \\
2014 & 47.4 & 18.0 & 24.8 & 43.6 & 11.4 & 14.472 \\
2015 & 48.2 & 17.9 & 24.5 & 43.5 & 10.6 & 10.687 \\
2016 & 47.5 & 19.2 & 23.7 & 44.3 & 9.2 & 15.874 \\
2017 & 47.3 & 18.5 & 24.6 & 44.7 & 10.6 & 14.967 \\
2018 & 48.5 & 17.9 & 24.5 & 43.7 & 9.7 & 14.323 \\
2019 & 49.0 & 17.3 & 24.5 & 43.1 & 9.8 & 13.591 \\
2020$^{\dagger}$ & 36.6 & 44.1 & 15.2 & 59.8 & 2.7 & 8.938 \\
2021$^{\dagger}$ & 41.9 & 32.0 & 21.1 & 53.7 & 6.1 & 12.181 \\
2022 & 35.7 & 43.7 & 15.7 & 59.7 & 6.6 & 12.497 \\
2023 & 37.4 & 42.7 & 14.7 & 57.7 & 5.8 & 12.351 \\
\bottomrule
\end{tabular}
\begin{tablenotes}\footnotesize
\item \textit{Fonte:} Painel longitudinal da PNADC trimestral (IBGE), 2012Q1--2024Q4.
\item \textit{Notas:} Promoção, repetência, evasão e não-progressão são taxas \emph{entre-anos} (entre o ano calendário $t$ e $t+1$). Abandono é \emph{intra-ano} (entre o início e o final de cada ano letivo).
\item $^{a}$Abandono refere-se ao percentual de alunos com \texttt{V3002=1} no 1º trimestre de $t$ que reportam não frequentar escola no 4º trimestre do mesmo ano. As taxas inter e intra não somam 100\% porque medem fenômenos em janelas temporais distintas.
\item $^{\dagger}$Anos afetados pela pandemia COVID-19 (suspensão de aulas presenciais, mudanças metodológicas na coleta da PNADC e adoção de aprovação automática em diversos estados). Interpretar com cautela.
\item \emph{N} reporta o número de observações individuais (não ponderado) usadas no cálculo. As taxas são ponderadas pelo peso amostral $w_i^{P1}$ (peso PNADC longitudinalmente válido).
\end{tablenotes}
\end{threeparttable}
\end{table}
